% =============================================================================
%  SECTION 7 — Conclusions & Future Work
%  Slides 7.1–7.3  |  Target time: ~5 minutes
%
%  Corresponds to dissertation Ch. 6 §§5–7
%
% =============================================================================

%%==========================================================================
%% CONCLUSIONS & FUTURE WORK
%%==========================================================================
\begin{frame}[shrink=15]{Novelty}
    \vspace{2.5em}

    \textbf{C1 — Policy Upsampling}
    \begin{itemize}\small
        \item First systematic evaluation of curriculum-based team scaling
            across \emph{heterogeneous} MARL environments
        \item Introduction of \textbf{agent-steps} as a normalized training
            cost metric for cross-configuration comparison
        \item Identification of observation schema design as the gating factor
            for policy transfer across team sizes
    \end{itemize}

    \textbf{C2 — Implicit Indication}
    \begin{itemize}\small
        \item First formalization of \textbf{implicit indication} — agent
            identity conveyed through observation structure, not explicit labels
        \item Theoretical guarantee (via Deep Sets) that a single shared policy
            is at least as expressive as the per-type collection
        \item Demonstrated zero-shot robustness to sensor dropout, team-size
            change, and novel compositions \emph{without retraining}
    \end{itemize}

    \textbf{C3 — Paradigm Comparison}
    \begin{itemize}\small
        \item First controlled comparison of architectural (GNN/PIC) vs.\
            representational invariance under matched conditions in HARL
        \item Finding that the simpler representational approach outperforms
            by $2.2$--$3.7\times$ across all configurations
    \end{itemize}

    \note{About 1.5 minutes.
    Walk through each contribution in one breath:
    ``C1 gave us the first systematic look at curriculum scaling
    in heterogeneous settings and a metric to measure training cost fairly.
    C2 gave us implicit indication — a single shared policy that works
    across structurally different agents, with a theoretical guarantee
    and emergent robustness we never trained for.
    C3 showed that this simple representational approach beats
    a purpose-built GNN architecture by two-to-four times.''
    Keep energy high — this is the `so what' for the committee.}
\end{frame}


\begin{frame}[shrink=7]{Minimum Tasks to Meet Objectives}
    \vspace{1em}
    \scriptsize
    % \renewcommand{\arraystretch}{1.15}
    \newcommand{\chit}{\item[\checkmark]}
    
    \begin{multicols}{2}
        \textbf{C1 — Policy Upsampling}
        \begin{itemize}
            \chit Design upsampling curriculum (pretrain $\to$ duplicate $\to$ retrain) % RT1.1
            \chit Define and validate agent-steps metric ($\rho > 0.999$) % RT1.2
            \chit Train tabula rasa baselines across all team sizes % RT1.3
            \chit Evaluate curriculum across pretraining lengths and team sizes % RT1.4
            \chit Select environments with distinct heterogeneity type % RT2.1
            \chit Adapt observation schemas for fixed-architecture transfer % RT2.2
            \chit Evaluate heterogeneity type effect on curriculum benefit % RT2.3
        \end{itemize}
        \vspace{0.4em}
        
        \textbf{C2 — Implicit Indication}
        \begin{itemize}
            \chit Design input-invariant (homogenized) architecture % RT1.1
            \chit Identify/build heterogeneous benchmark environment (HyperGrid) % RT1.2
            \chit Train and measure convergence across sensor configurations % RT1.3
            \chit Compare against HAPPO per-type baseline % RT1.4
            \chit Simulate sensor failure and degradation conditions % RT2.1
            \chit Simulate dynamic team-size changes and novel compositions % RT2.2
            \chit Measure stability and reward degradation under perturbation  % RT2.3
            \chit Statistical comparison across 8 perturbation conditions % RT2.4
            \chit Benchmark training/inference cost vs.\ per-type policies % RT3.1
            \chit Statistical comparison of convergence rate and degradation % RT3.2
        \end{itemize}
        \vspace{0.4em}
        
        \textbf{C3 — Paradigm Comparison}
        \begin{itemize}
            \chit Implement PIC under matched experimental conditions % RT1.1
            \chit Measure convergence and final performance vs.\ C2 baselines % RT1.2
            \chit Compare learning efficiency: PIC vs.\ implicit indication % RT1.3
            \chit Evaluate robustness to agent dropout and team-size changes % RT2.1
            \chit Random channel masking evaluation % RT2.2
            \chit Compare perturbation curves with C2 results % RT2.3
            \chit Benchmark computational cost of each paradigm % RT3.1
            \chit Statistical comparison: architectural vs.\ representational % RT3.2
        \end{itemize}
    \end{multicols}

    \note{
        This slide is presented to recall the research tasks discussed with 
        the committee in the prospectus.
    }
\end{frame}


\begin{frame}{Contribution Status}
\vspace{0.5em}
\begin{center}
\renewcommand{\arraystretch}{1.6}
\begin{tabular}{@{} l c l l @{}}
    \toprule
    \textbf{Contribution} & \textbf{Chapter} & \textbf{Status} & \textbf{Venue} \\
    \midrule
    C1: Policy Upsampling      & Ch.\ 3 & In Review & Autonomous Agents and Multi-Agent Systems \\
    C2: Implicit Indication    & Ch.\ 4 & In Review & Journal of Artificial Intelligence Research \\
    C3: Paradigm Comparison    & Ch.\ 5 & Awaiting Feedback& Committee \\
    \bottomrule
\end{tabular}
\end{center}

    \note{State this briefly and move on.
    The committee already knows the status — this slide is for the record.
    Simply say: ``All three contributions are complete and in the document you have in front of you.''
    If publication venues update before the defense, fill in the Venue column accordingly
    (e.g., journal name, conference acronym, or ``In Review'').
    Total time target: 30 seconds.}
\end{frame}


% # TODO: Consider consolidated Futures Slide:
% % -----------------------------------------------------------------------------
% % Slide 7.3 — Limitations & Future Directions
% % -----------------------------------------------------------------------------
% \begin{frame}{Limitations \& Future Directions}

%     \begin{columns}[t]
%         \begin{column}{0.45\textwidth}
%             \textbf{Limitations}
%             \begin{itemize}\small
%                 \item Single algorithm family (PPO-based);
%                     generalization to value-decomposition or
%                     off-policy methods is untested.
%                 \item Controlled simulation environments;
%                     maximum 8 agents — real swarms will be larger.
%                 \item Semantic decomposability assumption may not
%                     hold for all sensor modalities.
%                 \item Limited action-space heterogeneity —
%                     all agents share the same action space.
%             \end{itemize}
%         \end{column}
%         \begin{column}{0.52\textwidth}
%             \textbf{Future Work}
%             \begin{itemize}\small
%                 \item \textbf{Action-space heterogeneity:}
%                     extend homogenization to heterogeneous action
%                     spaces via action masking.
%                 \item \textbf{Dynamic team composition:}
%                     agents entering/leaving during training,
%                     not just at evaluation.
%                 \item \textbf{Hybrid approaches:}
%                     combine representational invariance with
%                     lightweight graph augmentation.
%                 \item \textbf{Physical validation:}
%                     noisy sensors, sim-to-real transfer,
%                     communication constraints.
%                 \item \textbf{Scale:}
%                     push toward Replicator-relevant swarm sizes
%                     ($50$--$100+$ agents).
%             \end{itemize}
%         \end{column}
%     \end{columns}

%     \vspace{0.3em}
%     \begin{center}
%         \small\emph{This dissertation opens more questions than it closes
%         — and that is a feature, not a bug.}
%     \end{center}

%     \note{About 1.5 minutes.
%     Be brief and honest on limitations —
%     the committee will respect candor more than defensiveness.
%     ``Four honest limitations: we used one algorithm family,
%     we stayed in simulation with small teams,
%     semantic decomposability is an assumption we did not relax,
%     and we did not vary action spaces.
%     Each of these is a natural next step.''
%     Then walk through the future work column:
%     ``Action-space heterogeneity is the most immediate extension —
%     homogenized action masking is a direct analogue of what C2 did
%     for observations.
%     Dynamic team composition during training, not just evaluation,
%     is the next frontier for robustness.
%     Hybrid representational-plus-architectural approaches could
%     capture relational structure that C3 showed pure GNNs miss.
%     And ultimately, physical validation and scaling toward
%     Replicator-relevant swarm sizes are where this work
%     connects back to the strategic motivation I opened with.''
%     End on the closing line — deliver it with a slight smile.
%     Then transition naturally to the closing slide:
%     ``With that, I'd like to thank my committee and open the floor
%     for questions.''}
% \end{frame}

